\documentclass[11pt,a4paper]{article}

% --- Packages (all part of a standard TeX Live / MiKTeX install) ----------
\usepackage[utf8]{inputenc}
\usepackage[T1]{fontenc}
\usepackage{amsmath}
\usepackage{amssymb}
\usepackage{booktabs}
\usepackage{graphicx}
\usepackage{siunitx}
\usepackage[hidelinks]{hyperref}

\sisetup{detect-weight=true, detect-family=true}

\title{Preprocessing Pipeline for Diabetic Retinopathy Grading\\
on the APTOS~2019 Fundus Image Dataset}
\author{AI for Health --- Group~5}
\date{\today}

\begin{document}
\maketitle

\begin{abstract}
We describe the preprocessing pipeline applied to the APTOS~2019 Blindness
Detection dataset prior to training a convolutional neural network for the
five-class grading of diabetic retinopathy (DR). The pipeline addresses four
characteristic challenges of colour fundus photography: heterogeneous image
quality and framing, low and locally varying lesion contrast, sensor noise, and
pronounced class imbalance. It comprises (i)~automated quality filtering to
discard truncated retinas, (ii)~retinal field-of-view cropping and resizing to
the input resolution of the target network, (iii)~mild edge-preserving noise
reduction, (iv)~green-channel extraction followed by Contrast Limited Adaptive
Histogram Equalisation (CLAHE) for contrast enhancement, (v)~normalisation of
pixel intensities to the unit interval, and (vi)~synthetic minority
over-sampling (SMOTE) applied exclusively to the training partition. The
resulting tensors are stored in a
model-ready format for an EfficientNet-B0 classifier. We report the exact
parameters and the per-class sample counts at each stage so that the procedure
is fully reproducible.
\end{abstract}

%==========================================================================
\section{Dataset}
\label{sec:dataset}

We use the APTOS~2019 Blindness Detection dataset, which contains colour
fundus photographs labelled by clinicians on the five-point International
Clinical Diabetic Retinopathy severity scale: \emph{0~--~no DR},
\emph{1~--~mild}, \emph{2~--~moderate}, \emph{3~--~severe}, and
\emph{4~--~proliferative DR}. The data are provided already partitioned into
training, validation, and test splits, each accompanied by a comma-separated
file mapping an image identifier (\texttt{id\_code}) to its grade
(\texttt{diagnosis}).

The images are acquired with a variety of fundus cameras under differing
conditions; consequently they are highly heterogeneous in resolution
(ranging from $480\times640$ to $2588\times3388$~pixels), aspect ratio, and
illumination. Many images are framed in landscape orientation, in which the
circular retinal field of view fills the full sensor height and is therefore
clipped at the top and bottom while the left and right margins remain black.
This is normal framing rather than image truncation, a distinction that
directly informs the quality-filtering criterion of
Section~\ref{sec:quality}. After verifying that the local and remotely hosted
copies of the dataset were byte-identical, the complete dataset
(\num{2930} training, \num{366} validation, and \num{366} test images) was
used as the source.

%==========================================================================
\section{Preprocessing Pipeline}
\label{sec:pipeline}

The pipeline is designed for an EfficientNet-B0 backbone, whose native input
is a $224\times224\times3$ tensor. The ordering of operations follows common
practice in the retinal-imaging literature: low-quality samples are discarded
first; geometric standardisation precedes intensity enhancement; and class
balancing is performed last, after the train/validation/test split, so that
synthetic samples can never leak into the evaluation partitions. The per-image
operations are summarised in Algorithm-style form below, followed by a
detailed description of each stage.

\begin{enumerate}
  \item \textbf{Quality filtering.} Detect and discard retinas that are
        genuinely truncated or over-cropped (Section~\ref{sec:quality}).
  \item \textbf{Field-of-view cropping.} Remove the uninformative black
        border surrounding the retinal disc (Section~\ref{sec:crop}).
  \item \textbf{Resizing.} Resample to $224\times224$~pixels
        (Section~\ref{sec:resize}).
  \item \textbf{Noise reduction.} Mild edge-preserving denoising to suppress
        sensor noise before contrast enhancement (Section~\ref{sec:denoise}).
  \item \textbf{Green-channel extraction and CLAHE.} Enhance vessel and
        lesion contrast (Section~\ref{sec:clahe}).
  \item \textbf{Normalisation.} Scale pixel intensities to the unit interval
        $[0,1]$ (Section~\ref{sec:normalise}).
  \item \textbf{Class balancing.} Apply SMOTE to the training partition only
        (Section~\ref{sec:smote}).
\end{enumerate}

%--------------------------------------------------------------------------
\subsection{Quality Filtering of Truncated Retinas}
\label{sec:quality}

A binary retina mask $M$ is obtained by converting the image to greyscale and
thresholding at an intensity of \num{15} (on the $[0,255]$ scale) to separate
the illuminated fundus from the dark background. The mask is regularised with
a morphological closing followed by an opening, both using a $7\times7$ square
structuring element, which fills small specular gaps and removes isolated
noise.

To decide whether a retina is truncated we examine four border bands, each of
width equal to \SI{1}{\percent} of the corresponding image dimension, located
at the left, right, top, and bottom edges. For each band we compute the mean
mask occupancy, i.e.\ the fraction of border pixels that belong to the retina.
A band is considered ``filled'' when its occupancy exceeds \num{0.30}.

A naive criterion that flags an image whenever \emph{any} single band is
filled is unsuitable for this dataset: the normally framed landscape images
of Section~\ref{sec:dataset} fill the top and bottom bands by construction,
so such a rule discards roughly \SI{60}{\percent} of the data and, critically,
the majority of the minority-grade samples. We therefore require that
\emph{at least three} of the four bands be filled, which is characteristic of
a retinal disc that has been clipped or over-cropped on multiple sides rather
than one that is merely framed in landscape. Formally, an image is discarded
when
\begin{equation}
  \sum_{s \in \{\text{L},\text{R},\text{T},\text{B}\}}
  \mathbb{1}\!\left[\,o_s > 0.30\,\right] \;\ge\; 3,
\end{equation}
where $o_s$ denotes the occupancy of border band $s$ and $\mathbb{1}[\cdot]$
is the indicator function. Under this criterion \num{81} of \num{2930}
training images (\SI{2.8}{\percent}) are removed, while the class distribution
is preserved almost exactly (Table~\ref{tab:counts}). The same filter is
applied to the validation and test partitions, removing \num{10} and \num{7}
images respectively.

%--------------------------------------------------------------------------
\subsection{Field-of-View Cropping}
\label{sec:crop}

Fundus photographs contain large black margins that carry no diagnostic
information and waste effective resolution after resizing. We crop each image
to the tight axis-aligned bounding box of its retinal region. The retinal
region is identified by thresholding the greyscale image at an intensity of
\num{20}; the bounding box is the smallest rectangle enclosing all pixels
above this threshold. If no pixel exceeds the threshold the image is returned
unchanged.

%--------------------------------------------------------------------------
\subsection{Resizing}
\label{sec:resize}

The cropped retina is resampled to a fixed $224\times224$ resolution to match
the input dimensions of EfficientNet-B0. Resizing uses area interpolation
(OpenCV's \texttt{INTER\_AREA}), which is well suited to the downscaling that
dominates this heterogeneous-resolution dataset.

%--------------------------------------------------------------------------
\subsection{Noise Reduction}
\label{sec:denoise}

Colour fundus photographs carry sensor and compression noise. Because the
subsequent CLAHE step (Section~\ref{sec:clahe}) amplifies local contrast and
would otherwise accentuate this noise, we denoise \emph{before} contrast
enhancement. We apply a mild edge-preserving bilateral filter (diameter
$d=5$, $\sigma_{\text{colour}}=\sigma_{\text{space}}=\num{35}$), which smooths
homogeneous regions while preserving the edges of vessels and lesions, the
diagnostically relevant structures. The conservative setting avoids erasing
small lesions such as microaneurysms. Alternative methods (non-local means,
median, Gaussian) are selectable behind a single parameter, with the bilateral
filter as the default. Denoising operates on the 8-bit image, before the
normalisation of Section~\ref{sec:normalise}.

%--------------------------------------------------------------------------
\subsection{Green-Channel Extraction and CLAHE}
\label{sec:clahe}

In RGB fundus imagery the green channel exhibits the highest contrast between
the retinal vasculature and lesions on the one hand and the background on the
other, while the red channel is frequently saturated and the blue channel is
noisy. We therefore retain only the green channel for subsequent processing.

Because illumination varies markedly across a single retina, we apply Contrast
Limited Adaptive Histogram Equalisation (CLAHE) to the green channel. CLAHE
equalises the histogram within local tiles rather than globally, and clips the
contribution of each histogram bin to limit noise amplification in homogeneous
regions. We use a clip limit of \num{2.0} and a tile grid of $8\times8$,
yielding locally enhanced contrast of microaneurysms, haemorrhages, and
exudates without introducing strong artefacts.

EfficientNet-B0 expects a three-channel input. As the informative signal is
now contained in a single enhanced channel, we replicate the CLAHE-processed
green channel across all three channels, producing a $224\times224\times3$
tensor whose channels are identical. This preserves compatibility with the
network stem and with ImageNet-pretrained weights while avoiding the
introduction of redundant or noisy colour information.

%--------------------------------------------------------------------------
\subsection{Intensity Normalisation}
\label{sec:normalise}

Finally, the pixel intensities are normalised from the $[0,255]$ integer range
to the $[0,1]$ unit interval by dividing by~\num{255},
\begin{equation}
  p_{\text{norm}} = \frac{p}{255},
\end{equation}
where $p$ is the original pixel value and $p_{\text{norm}}$ its normalised
counterpart. Normalising the dynamic range across images acquired under
different illumination standardises the input statistics and stabilises and
accelerates the optimisation of the downstream network. The arrays are
consequently stored in single-precision floating point
(Section~\ref{sec:output}).

A practical caveat for deployment: the Keras \texttt{EfficientNetB0} model
provides its own rescaling layer through
\texttt{efficientnet.preprocess\_input}. Because the data produced here is
already scaled to $[0,1]$, it must be fed to the network \emph{without} a
second application of \texttt{preprocess\_input}, otherwise the scaling would
be applied twice.

%--------------------------------------------------------------------------
\subsection{Class Balancing with SMOTE}
\label{sec:smote}

The training partition is strongly imbalanced, with the \emph{no~DR} grade
accounting for roughly half of all samples and the \emph{severe} grade for
only about \SI{5}{\percent} (Table~\ref{tab:counts}). To prevent the
classifier from being biased towards the majority grade we balance the
training set with the Synthetic Minority Over-sampling Technique (SMOTE).
Each preprocessed image is flattened into a feature vector of
$224\times224\times3 = \num{150528}$ values, and SMOTE synthesises new
minority-class samples by interpolating between a sample and its
$k$~nearest neighbours of the same class ($k=5$). After resampling, the
vectors are reshaped back to image tensors.

Two points are essential for correctness. First, SMOTE is fitted on the
\emph{training partition only}; the validation and test partitions retain
their natural class distribution so that performance is estimated on a
realistic prior. Second, SMOTE operates on the already-normalised floating
point images of Section~\ref{sec:normalise}, which sidesteps a subtle failure
mode: performed directly on unsigned 8-bit pixels, the interpolation
$\mathbf{x}_{\text{new}} = \mathbf{x}_i + \lambda\,(\mathbf{x}_j - \mathbf{x}_i)$,
$\lambda \in [0,1]$, would wrap around modulo~256 wherever the difference
$\mathbf{x}_j - \mathbf{x}_i$ is negative and yield corrupted synthetic images.
Because both operands already lie in $[0,1]$, every interpolated sample also
lies in $[0,1]$, so no rounding or clipping is required. The feature matrix is
kept in single-precision floating point, which halves the peak memory footprint
relative to the double precision SMOTE would otherwise use.

Balancing raises every training grade to the majority count of \num{1411},
giving a balanced training set of \num{7055} images.

\begin{table}[t]
  \centering
  \caption{Per-class image counts at each stage of the pipeline. The
  training set is shown before quality filtering, after filtering, and after
  SMOTE balancing; the validation and test sets are shown after filtering
  (no balancing is applied to them). Grades follow the five-point clinical
  DR scale (0--4).}
  \label{tab:counts}
  \begin{tabular}{lrrrrrr}
    \toprule
    & \multicolumn{5}{c}{DR grade} & \\
    \cmidrule(lr){2-6}
    Stage & 0 & 1 & 2 & 3 & 4 & Total \\
    \midrule
    Train (raw)            & 1434 & 300 & 808 & 154 & 234 & 2930 \\
    Train (after filter)   & 1411 & 294 & 770 & 145 & 229 & 2849 \\
    Train (after SMOTE)    & 1411 & 1411 & 1411 & 1411 & 1411 & 7055 \\
    \midrule
    Validation (after filter) & 169 & 40 & 98 & 21 & 28 & 356 \\
    Test (after filter)       & 196 & 29 & 84 & 17 & 33 & 359 \\
    \bottomrule
  \end{tabular}
\end{table}

%==========================================================================
\section{Output Format}
\label{sec:output}

Each partition is stored as a pair of arrays: an image tensor of shape
$(N, 224, 224, 3)$ in single-precision floating point with values in $[0,1]$
and a corresponding integer label vector of length $N$. As noted in
Section~\ref{sec:normalise}, because the data is already normalised it should
be passed to the Keras EfficientNet-B0 model \emph{without} a further call to
its \texttt{preprocess\_input} rescaling layer, which would otherwise apply the
scaling twice. The synthetic samples produced by SMOTE have no natural file
identifier, so the partitions are stored as serialised arrays rather than as
individual image files.

%==========================================================================
\section{Reproducibility}
\label{sec:repro}

Table~\ref{tab:params} lists every parameter required to reproduce the
pipeline. All randomised operations (SMOTE) use a fixed seed of \num{42}.

\begin{table}[t]
  \centering
  \caption{Pipeline parameters.}
  \label{tab:params}
  \begin{tabular}{lll}
    \toprule
    Stage & Parameter & Value \\
    \midrule
    Retina mask        & Greyscale threshold        & 15 \\
                       & Morphological kernel       & $7\times7$ \\
                       & Operations                 & close, then open \\
    \midrule
    Quality filter     & Border-band width          & \SI{1}{\percent} of side \\
                       & Occupancy threshold        & 0.30 \\
                       & Filled-bands required      & $\ge 3$ of 4 \\
    \midrule
    FOV crop           & Greyscale threshold        & 20 \\
    \midrule
    Resize             & Target resolution          & $224\times224$ \\
                       & Interpolation              & area (\texttt{INTER\_AREA}) \\
    \midrule
    Noise reduction    & Filter (default)           & bilateral \\
                       & Diameter $d$               & 5 \\
                       & $\sigma_{\text{colour}}$, $\sigma_{\text{space}}$ & 35 \\
    \midrule
    Contrast           & Channel                    & green \\
                       & CLAHE clip limit           & 2.0 \\
                       & CLAHE tile grid            & $8\times8$ \\
                       & Output channels            & green replicated $\times 3$ \\
    \midrule
    Normalisation      & Scaling                    & divide by 255 \\
                       & Output range               & $[0,1]$ \\
                       & Output dtype               & float32 \\
    \midrule
    Balancing          & Method                     & SMOTE \\
                       & Neighbours $k$             & 5 \\
                       & Working precision          & float32 \\
                       & Random seed                & 42 \\
                       & Applied to                 & training set only \\
    \bottomrule
  \end{tabular}
\end{table}

%==========================================================================
\begin{thebibliography}{9}

\bibitem{aptos2019}
Asia Pacific Tele-Ophthalmology Society,
\emph{APTOS 2019 Blindness Detection},
Kaggle, 2019.

\bibitem{efficientnet}
M.~Tan and Q.~V.~Le,
``EfficientNet: Rethinking Model Scaling for Convolutional Neural Networks,''
\emph{Proc.\ Int.\ Conf.\ on Machine Learning (ICML)}, 2019.

\bibitem{clahe}
S.~M.~Pizer \emph{et al.},
``Adaptive Histogram Equalization and Its Variations,''
\emph{Computer Vision, Graphics, and Image Processing}, vol.~39, no.~3,
pp.~355--368, 1987.

\bibitem{smote}
N.~V.~Chawla, K.~W.~Bowyer, L.~O.~Hall, and W.~P.~Kegelmeyer,
``SMOTE: Synthetic Minority Over-sampling Technique,''
\emph{Journal of Artificial Intelligence Research}, vol.~16, pp.~321--357,
2002.

\end{thebibliography}

\end{document}
